% !PRINTED 436
\weberpage{436}
%<<CONT
sieht auch leicht, dass diese $n$ Werthe alle von einander verschieden sind, wenn man von dem besonderen Falle absieht, dass $\varphi$ ein Vielfaches von $\pi$ ist. Denn zwei Cosinus sind nur dann einander gleich, wenn die Summe oder die Differenz der Winkel ein Vielfaches von $2\pi$ ist.

Die Gleichung
\weq{15}{\sin n\varphi = \sin\varphi\, B_n(x)}
bestimmt, wenn $B_n(x)$ nicht gleich Null ist, $\sin\varphi$ eindeutig durch $x$.

Suchen wir in der Gleichung (13) das von $x$ unabhängige Glied, so ergiebt sich mit Benutzung von (6) für ein gerades $n$
\[
2(-1)^{\frac{n}{2}} - 2\cos n\varphi,
\]
und für ein ungerades $n$
\[
-2\cos n\varphi,
\]
und wenn man dies dem positiven oder negativen Product der Wurzeln gleichsetzt, so erhält man für ein gerades $n$
\weq{16}{2^{n-1} \overset{\nu}{\Pi} \cos\left(\varphi + \dfrac{2\nu\pi}{n}\right) = (-1)^{\frac{n}{2}} - \cos n\varphi,}
und für ein ungerades $n$
\weq{17}{2^{n-1} \overset{\nu}{\Pi} \cos\left(\varphi + \dfrac{2\nu\pi}{n}\right) = \cos n\varphi,}
und diese Formeln gelten, da rechts und links stetige Functionen von $\varphi$ stehen, auch noch für den zunächst ausgeschlossenen Fall, wo $\varphi$ ein Vielfaches von $\pi$ ist.

Hierin können wir unter dem Productzeichen $\nu m$ für $\nu$ setzen, wenn $m$ irgend eine zu $n$ theilerfremde Zahl ist, weil dann $\nu m$ zugleich mit $\nu$ ein volles Restsystem nach dem Modul $n$ durchläuft. Sondern wir den Werth $\nu=0$ ab, und lassen dann $\nu$ sowohl die positiven als die negativen Zahlen durchlaufen, die absolut kleiner als $n:2$ sind, so erhalten wir aus (17) die für jedes ungerade $n$ gültige Formel
\weq{18}{2^{n-1} \overset{\nu}{\underset{1,\,\frac{n-1}{2}}{\Pi}} \cos\left(\varphi + \dfrac{2\nu m\pi}{n}\right) \cos\left(\varphi - \dfrac{2\nu m\pi}{n}\right) = \dfrac{\cos n\varphi}{\cos\varphi}.}
Setzen wir hierin $\varphi = 0$ und ziehen die Quadratwurzel, so ergiebt sich
%CONT>>
